\chapter{Final Cipher Selection}

\section{Decision}

The final cipher selected for the authenticated encryption pipeline is
AES-128 in CTR (Counter) mode for encryption. CTR mode converts the AES block
cipher into a stream cipher: the IV (initialisation vector) and a counter are
concatenated to form a block input, AES encrypts that block to produce a
keystream block, and the keystream is XORed with the plaintext to produce
ciphertext. Crucially, CTR mode requires only the forward AES block cipher
(not the inverse), and encryption and decryption are the same operation. The
IV must be unique for each message encrypted under the same key; reusing an
IV with the same key would allow an adversary to recover plaintext by XORing
two ciphertexts together.

Authentication is provided by HMAC-SHA256 (truncated to 192 bits), appended
to the AES-CTR ciphertext in an Encrypt-then-MAC pattern. HMAC (Hash-based
Message Authentication Code) computes a keyed hash over the authenticated
data using a cryptographic hash function --- in this case SHA-256. HMAC
provides integrity and authenticity: an adversary who does not know the MAC
key cannot produce a valid tag for a tampered ciphertext.

At implementation level, this decision is realized as two explicit repository
targets: \texttt{aes-ctr} for confidentiality-only benchmarking and
\texttt{aes-ctr-hmac} for authenticated Encrypt-then-MAC benchmarking. The split
is intentional and supports clearer measurement and traceability.

\section{Justification}

The decision is based on measured behavior in this zkVM environment. AES-128
provides better empirical proving performance than LowMC in the tested
implementations, and the AES-CTR path integrates cleanly with the project's
Encrypt-then-MAC construction under FIPS 197 and NIST SP 800-38A
\cite{nistfips197,nistsp80038a}.

This is an explicitly practical choice. Although AES is not usually presented as
the most ``ZK-specialized'' cipher in theory discussions, the project
prioritizes measured zkVM behavior over abstract circuit expectations.

The four decision criteria --- performance, security maturity, implementation
risk, and proof-integration clarity --- were applied as follows.
Table~\ref{tab:auth-alternatives} in Chapter 7 shows the systematic comparison
of AES-CTR+HMAC against the main alternatives (AES-GCM, ChaCha20-Poly1305, and
OCB) using these same criteria. Performance determines whether the final pipeline
is feasible to prove. Security maturity determines whether the primitive and mode
can be justified with standard assumptions. Implementation risk captures how much
new unaudited code must be added. Proof-integration clarity captures whether the
resulting transcript can be explained cleanly in the proof-output IND-CPA game.

\section{AES Backend Comparison Snapshot}

The remaining engineering question is whether the AES optimization workspace
improves proving performance under strict like-for-like settings.

A direct snapshot from the baseline comparison campaign gives:

\begin{table}[t]
\centering
\small
\begin{tabular}{lrrr}
\toprule
Target & Prove s (median) & User cycles (median) & Delta vs \texttt{aes-r0} \\
\midrule
\texttt{aes-r0} & 14.8429 & 181710 & baseline \\
\texttt{aes-r0-optimised} & 14.8359 & 181710 & -0.05\% prove, 0.00\% cycles \\
\bottomrule
\end{tabular}
\caption{AES baseline versus optimisation workspace (single-trial, same run).}
\label{tab:aes-backend-snapshot}
\end{table}

Table~\ref{tab:aes-backend-snapshot} shows no meaningful measured improvement
from the S-box backend change in that single-trial run. Therefore the final
cipher decision is anchored to stable baseline behavior and integration quality,
not to a claimed backend optimization advantage.

Regarding the HMAC hash function: this project benchmarks AES-CTR as the
stream cipher component; benchmarking the SHA-256 hash function used inside
HMAC as a standalone zkVM target is out of scope for the current evaluation
campaign. The HMAC overhead is captured indirectly in the CTR+HMAC versus
CTR-only comparison in Chapter 11, where the user-cycle delta reflects the full
in-guest HMAC computation including SHA-256 evaluation. A dedicated SHA-256
benchmark target would separate this cost further and is left as future work.

A stricter AES backend study (explicit \texttt{sbox-table} vs
\texttt{sbox-logic} feature flags under repeated trials) also remains future
work.

\section{Decision Summary}

The selected option is not the fastest possible authenticated-encryption design
in general; it is the most defensible design under this project's evidence. AES
beats the available LowMC implementation by a large measured margin, CTR mode
keeps encryption behavior simple to benchmark, and HMAC provides a separate
authentication layer whose overhead can be measured directly. The trade-off is
that the final construction is two-pass and therefore less elegant than an
integrated AEAD mode such as GCM or OCB. That cost is accepted because the
separation improves artifact traceability and simplifies the security argument.
